\chapter{John F. Nash Jr. \& Louis Nirenberg (2015)}

\section{Biographical Sketches}



\subsection{John Forbes Nash Jr.}

John Nash was born on 13 June 1928 in Bluefield, West Virginia, USA. He studied at Carnegie Institute of Technology (now Carnegie Mellon University) and Princeton University, where he completed his PhD in 1950 under the supervision of Albert W. Tucker. His 27-page doctoral dissertation introduced the concept of the Nash equilibrium, which revolutionized economics. Nash held positions at the Massachusetts Institute of Technology and Princeton University. He received the Nobel Prize in Economics in 1994 and the Abel Prize in 2015. Nash's life and struggle with schizophrenia were the subject of the book and film "A Beautiful Mind."

\subsection{Louis Nirenberg}

Louis Nirenberg was born on 28 February 1925 in Hamilton, Ontario, Canada. He studied at McGill University and New York University, where he completed his PhD in 1949 under the supervision of James Serrin and Kurt Friedrichs. He spent his entire career at the Courant Institute of Mathematical Sciences at NYU. Nirenberg received the Abel Prize (jointly with Nash) in 2015, the National Medal of Science (1995), and the Chern Medal (2010).

\section{High-Level View for a General Audience}

\subsection{What Did Nash and Nirenberg Do?}

Think about the shape of a soap film stretched across a wire frame. The film takes the shape that minimizes its surface area. This is a problem of geometry: given a boundary, find the surface of least area spanning it.

Now imagine that you are told the curvature of a surface at every point, but not its overall shape. Can you reconstruct the shape from this curvature information? This is the kind of problem that John Nash and Louis Nirenberg solved.

\paragraph{Nash's Contribution:} Nash proved that any Riemannian manifold can be isometrically embedded into Euclidean space with sufficiently many dimensions. In simple terms: any curved space can be placed inside a flat space of higher dimension in a way that preserves all distances. This is the Nash embedding theorem, one of the most surprising results in geometry.

\paragraph{Nirenberg's Contribution:} Nirenberg made fundamental contributions to the theory of partial differential equations, particularly equations describing surfaces and their curvatures. He developed the tools needed to prove the existence and regularity of solutions to these geometric PDEs.

\subsection{Why Does It Matter?}

\begin{itemize}
\item \textbf{Geometry and Physics}: The Nash embedding theorem tells us that curved spaces (like spacetime in general relativity) can be studied as subsets of flat spaces of higher dimension. This is a key insight for both mathematics and physics.

\item \textbf{Materials Science}: The mathematics of surfaces and their curvatures, which Nash and Nirenberg studied, is used to understand the behavior of thin films, crystals, and other materials.

\item \textbf{Image Processing}: The partial differential equations that Nirenberg studied are used in image processing, including edge detection, denoising, and image segmentation.

\item \textbf{Animation and Graphics}: The reconstruction of surfaces from geometric data, which uses the mathematics of Nash and Nirenberg, is used in computer graphics and animation.

\item \textbf{Economics}: While separately famous for Nash equilibrium in game theory, Nash's mathematical work on embedding theorems has deep connections to the geometry of economic models and their representations.
\end{itemize}

\section{The Origin of the Problem}

\subsection{The Embedding Problem}

The embedding problem asks: given an abstract Riemannian manifold $(M^n, g)$, can it be realized as a submanifold of Euclidean space $\mathbb{R}^N$ in such a way that the induced metric equals $g$?

The isometric embedding problem was posed by Riemann himself in his 1854 habilitation lecture. By the 1950s, the following was known:
\begin{itemize}
\item Local isometric embeddings exist for analytic metrics (Cauchy--Kovalevskaya theorem).
\item The Whitney embedding theorem provides topological embeddings into $\mathbb{R}^{2n}$.
\item But isometric embeddings---those that preserve distances---seemed much harder.
\end{itemize}

\subsection{The Geometric PDE Problem}

Many problems in differential geometry reduce to solving partial differential equations. For example, the prescribed curvature problem asks: given a function $K$ on a surface, can we find a metric whose Gaussian curvature is $K$? This leads to the Yamabe equation and the Monge--Amp\`ere equation.

These equations are \emph{nonlinear} and \emph{elliptic}. Their analysis requires powerful techniques that Nash and Nirenberg helped develop.

\section{Deep Insight into the Problem}

\subsection{Nash's $C^1$ Embedding Theorem}

In 1954, Nash stunned the mathematical world by proving that any Riemannian manifold can be isometrically embedded into Euclidean space. His first proof (the $C^1$ embedding theorem) used a remarkably simple idea: a series of small perturbations that gradually reduce the error in the embedding.

The $C^1$ Nash embedding theorem: every compact Riemannian $n$-manifold admits a $C^1$ isometric embedding into $\mathbb{R}^{n+1}$.

The construction works as follows:
\begin{enumerate}
\item Start with a short embedding (one that shrinks all distances).
\item Decompose the difference between the metric $g$ and the pullback of the Euclidean metric into a sum of small "elementary" metrics.
\item Each elementary metric corresponds to a perturbation that adds a wrinkle to the embedding in a new direction.
\item With infinitely many wrinkles in $n+1$ directions, the metric error is corrected without changing the $C^1$ topology.
\end{enumerate}

\subsection{The Smooth Nash Embedding Theorem}

The smooth version requires more dimensions. Nash proved:
\[
(M^n, g) \text{ embeds isometrically into } \mathbb{R}^{n(3n+11)/2}.
\]

The dimension bound was later improved by Gromov and Gunther to $\frac{n(n+1)}{2} + 2n + 3$ and then to $\max\{ n(n+1)/2 + n, n(n+1)/2 + 5 \}$.

\subsection{The Nash--Moser Implicit Function Theorem}

Nash's proof of the smooth embedding theorem required a new kind of implicit function theorem, now called the Nash--Moser theorem. Unlike the classical inverse function theorem, which works in Banach spaces, the Nash--Moser theorem works in Fr\'echet spaces and requires smoothing operations at each step of the iteration.

The technical innovation: the iteration scheme includes a smoothing operator $S_\theta$ that regularizes the solution at each step. The "tame" estimate ensures that the smoothing does not amplify errors. The theorem was refined by Moser (hence the name Nash--Moser) and later by Gromov and Hamilton.

\subsection{The Nash Twist}

The key geometric construction in Nash's proof is the "Nash twist"---a periodic deformation of the embedding that adds a short-wavelength correction to the metric without creating self-intersections. The twist is controlled by a rapidly oscillating function, and the product structure ensures that the different corrections do not interfere.

\subsection{Nirenberg's Work on Elliptic PDEs}

Nirenberg made fundamental contributions to the theory of linear and nonlinear elliptic PDEs:
\begin{itemize}
\item \textbf{Schauder Estimates}: Nirenberg developed a priori estimates for solutions of elliptic equations of arbitrary order, including the interior Schauder estimates and the boundary estimates. These estimates control the $C^{k,\alpha}$ norm of the solution in terms of the $C^{k-2,\alpha}$ norm of the equation.
\item \textbf{Aglmon--Douglis--Nirenberg (ADN) Theory}: With Agmon and Douglis, Nirenberg developed a general theory of elliptic boundary value problems for systems of equations. The ADN theory is the standard framework for analyzing elliptic boundary value problems.
\item \textbf{The Method of Moving Planes}: Nirenberg developed this technique for proving symmetry of solutions to nonlinear PDEs. The idea: reflect the solution across a plane and compare the original with the reflection using the maximum principle. If the solution is monotone up to a critical plane, it must be symmetric.
\end{itemize}

\subsection{The Gagliardo--Nirenberg--Sobolev Inequality}

This fundamental inequality controls the $L^{p^*}$ norm of a function in terms of its gradient:
\[
\|u\|_{L^{p^*}(\mathbb{R}^n)} \leq C \|\nabla u\|_{L^p(\mathbb{R}^n)},
\]
where $p^* = \frac{np}{n-p}$ is the Sobolev conjugate exponent. This inequality is essential for establishing compactness in PDE theory.

\section{Biggest Contributions and New Tools}

\subsection{The Nash Embedding Theorem}

Nash's embedding theorem comes in two versions:
\[
\text{Every smooth Riemannian manifold } (M^n, g) \text{ embeds isometrically into } \mathbb{R}^{n(3n+11)/2}.
\]

The proof introduced the "Nash twist" and the iteration method that came to be called Nash--Moser iteration.

\subsection{The Nash--Moser Theorem}

This theorem provides an inverse function theorem in Fr\'echet spaces under the conditions of a "tame" estimate and a smooth approximation scheme. The theorem is essential for Nash's smooth embedding theorem and has applications in KAM theory and the geometry of diffeomorphism groups.

\subsection{Nirenberg's Maximum Principle and the Moving Planes Method}

Nirenberg's maximum principle for elliptic equations states that the maximum of a solution occurs on the boundary unless the solution is constant. This simple but powerful tool is used throughout PDE theory.

\subsection{The Gagliardo--Nirenberg Interpolation Inequality}

This inequality estimates intermediate derivatives in terms of lower and higher derivatives:
\[
\|D^j u\|_{L^p} \leq C \|D^m u\|_{L^r}^\theta \|u\|_{L^q}^{1-\theta}.
\]

\subsection{The Brezis--Nirenberg Problem}

The Brezis--Nirenberg problem concerns the existence of positive solutions to the critical Sobolev equation:
\[
-\Delta u = u^{(n+2)/(n-2)} + \lambda u \quad \text{in } \Omega, \quad u = 0 \text{ on } \partial\Omega.
\]

\section{Impact of the Discovery}

\subsection{Impact on Mathematics}

\begin{itemize}
\item \textbf{Differential Geometry}: The Nash embedding theorem is a fundamental result that justifies the study of Riemannian manifolds as subsets of Euclidean space.

\item \textbf{PDE Theory}: Nirenberg's work on elliptic PDEs provides essential tools for geometric analysis, the theory of minimal surfaces, and the study of complex Monge--Amp\`ere equations.

\item \textbf{Functional Analysis}: The Nash--Moser theorem is a powerful tool in the analysis of nonlinear PDEs in geometry.

\end{itemize}

\subsection{Impact on Technology}

\begin{itemize}
\item \textbf{Computer Vision}: The level set method and the Perona--Malik equation, based on PDEs of the type studied by Nirenberg, are used in image processing.

\item \textbf{Materials Science}: The geometry of thin films and interfaces uses the theory of surfaces and curvature that Nash and Nirenberg advanced.

\end{itemize}

\subsection{Impact on Economics}

Nash's equilibrium concept, now called "Nash equilibrium," transformed economics and social science. It provides the mathematical foundation for analyzing strategic interactions in markets, politics, and social behavior.

\section{Current Developments}

\subsection{Isometric Embedding in Low Codimensions}

The minimal dimension needed for isometric embeddings remains an active area of research. Gunther improved Nash's bound to $N = \max\{n(n+1)/2 + n, n(n+1)/2 + 5\}$. The optimal dimension is conjectured to be $n(n+1)/2$ for compact manifolds and $n(n+1)/2 + n$ for non-compact ones.

\subsection{Fluid Dynamics and the Nash--Moser Theorem}

The Nash--Moser theorem has been applied to the study of the Euler equations. In particular, De Lellis and Sz\'ekelyhidi used Nash-type iteration to construct non-unique weak solutions of the Euler equations, establishing the failure of Onsager's conjecture for low regularity.

\subsection{Nonlinear PDEs in Geometry}

Nirenberg's methods continue to be essential in the study of:
\begin{itemize}
\item \textbf{The Yamabe Problem}: Finding a metric of constant scalar curvature.
\item \textbf{The Prescribed Scalar Curvature Problem}: Given a function $K$, does there exist a metric with scalar curvature $K$?
\item \textbf{The Ricci Flow}: The parabolic version of the Yamabe problem, crucial for Perelman's proof of the Poincar\'e conjecture.
\end{itemize}

\subsection{Game Theory and Machine Learning}

Nash equilibrium concepts are used in multi-agent reinforcement learning and the design of machine learning algorithms. The theory of generative adversarial networks (GANs) is based on a two-player game, and the training dynamics are analyzed using game theory.

\subsection{The Isometric Embedding of the Sphere}

The isometric embedding of the sphere $(S^2, g_0)$ into $\mathbb{R}^3$ is a classical problem. While $S^2$ cannot be isometrically embedded into $\mathbb{R}^2$, the $C^1$ embedding theorem says it can be embedded into $\mathbb{R}^3$ with a $C^1$ but not $C^2$ map (this is the subject of the Nash--Kuiper theorem, which constructs $C^1$ isometric embeddings of surfaces into $\mathbb{R}^3$ that are arbitrarily close to any short embedding).

\section{Conclusion}

John Nash and Louis Nirenberg transformed our understanding of differential geometry and partial differential equations. Nash's embedding theorem is one of the most surprising results in geometry, showing that any curved space can be realized as a subset of flat space. Nirenberg's work on PDEs provided the tools needed to solve the geometric equations that arise in the study of surfaces and their curvatures.

Together, they received the Abel Prize "for striking and seminal contributions to the theory of nonlinear partial differential equations and its applications to geometric analysis."


\section{Behind the Breakthrough: The Human Story}
Nash's embedding theorem was born from a bet. A colleague claimed Nash couldn't prove that any Riemannian manifold could be embedded in Euclidean space. Nash went to work and proved it, proving the colleague wrong.

\section{Pedagogical Metaphors and Lineage}
\subsection{Signature Metaphor}
``Folding and weaving any curved space smoothly into Euclidean space without stretching it, like wrapping a gift without creasing the paper.''

\subsection{Mathematical Lineage}
Nash advised by Albert W. Tucker. Nirenberg advised by James Stoker.

\section{Applied Science and Computational Algorithms}
\subsection{Key Takeaways for Applied Science}
Applied in computer graphics (isometric texture mapping, mesh deformation, and skinning) and in aerospace engineering to solve fluid flow and heat dissipation boundaries.

\subsection{Algorithms and Numerical Implementations}
Least Squares Conformal Maps (LSCM) and geometric flattening algorithms in CAD software solve boundary value problems using Nash's embedding coordinates.

\section{The Research Frontier}
Proving global regularity or blow-up of solutions to the 3D Navier-Stokes equations, which is a major open problem in mathematical physics.

\section{Guided Exercises and Challenges}
\begin{enumerate}[label=\textbf{\arabic*.}]
  \item State the Nash embedding theorem and explain the difference between $C^1$ isometric embeddings and smooth $C^k$ embeddings.
  \item Explain the Nash-Moser inverse function theorem and why it is needed when solving non-linear PDEs where derivative loss occurs.
  \item State the De Giorgi-Nash-Moser theorem on the H\"older continuity of weak solutions to elliptic equations.
\end{enumerate}

\section{Curriculum Map and Further Reading}
For students wishing to delve deeper into the mathematical machinery underlying these breakthroughs, the following learning path is recommended:
\begin{itemize}
  \item John Nash, \emph{The Imbedding Problem for Riemannian Manifolds}, Annals of Mathematics 63 (1956).
  \item L. Nirenberg, \emph{Topics in Nonlinear Functional Analysis}, AMS.
\end{itemize}

\section*{Bibliography}
\begin{enumerate}
\item J. F. Nash Jr., "$C^1$ isometric imbeddings," Annals of Mathematics 60 (1954), 383--396.
\item J. F. Nash Jr., "The imbedding problem for Riemannian manifolds," Annals of Mathematics 63 (1956), 20--63.
\item J. F. Nash Jr., "Non-cooperative games," Annals of Mathematics 54 (1951), 286--295.
\item J. F. Nash Jr., "Equilibrium points in $n$-person games," Proceedings of the National Academy of Sciences 36 (1950), 48--49.
\item L. Nirenberg, "Topics in Nonlinear Functional Analysis," New York University, 1974.
\item S. Agmon, A. Douglis, and L. Nirenberg, "Estimates near the boundary for solutions of elliptic partial differential equations," Communications on Pure and Applied Mathematics 12 (1959), 623--727.
\item L. Nirenberg, "The Weyl and Minkowski problems in differential geometry," Communications on Pure and Applied Mathematics 6 (1953), 337--394.
\item L. Nirenberg, "On elliptic partial differential equations," Annali della Scuola Normale Superiore di Pisa 13 (1959), 115--162.
\item H. Brezis and L. Nirenberg, "Positive solutions of nonlinear elliptic equations involving critical Sobolev exponents," Communications on Pure and Applied Mathematics 36 (1983), 437--477.
\item R. S. Hamilton, "The inverse function theorem of Nash and Moser," Bulletin of the American Mathematical Society 7 (1982), 65--222.
\end{enumerate}
